\begin{musica}{Cielito Lindo}{Popular Mexicana}{Tuna de Engenharia da Universidade do Porto}
\textbf{[}\ac{A}{Ese} lu\ac{E7}{nar} que \ac{A}{tie}nes, cie\ac{E7}{li}to lin\ac{A}{do}, junto \ac{F#7}{a} la \ac{Bm}{bo}ca
\ac{E7}{No} te \ac{Bm}{lo} des a \ac{E7}{na}die, \ac{Bm}{cie}lito \ac{E7}{lin}do, que a mi me \ac{A}{to}ca.\textbf{]} \nota{2x}

\nota{Refrão:}
\ac{A}{Ay}, \ac{A7}{ay}, ay, \ac{D}{ay} \ac{F#}{}__
\ac{Bm}{Can}ta \ac{E7}{y} no \ac{A}{llo}res
Porque \ac{F#7}{can}tando se a\ac{Bm}{le}gra, \ac{E7}{cie}lito lindo, los cora\ac{A}{zo}nes.

De la sierra morena, cielito lindo, vienen bajando
Un par de ojitos niegros, cielito lindo, de contrabando.

Una flecha en el aire, cielito lindo, tiró Cupido
Él la tir'o julgando, cielito lindo, y a mi me toca.
\nota{D: xx0232}
\nota{E7: 020100}
\nota{F#: 244322}
\nota{F#7: 242322}
\nota{A: x02220}
\nota{A7: x02020}
\nota{Bm: x24432}

\posicoes{D: xx0232, E7: 020100, F#: 244322, F#7: 242322, A: x02220, A7: x02020,
  Bm: x24432}
\end{musica}
